\section{结构性综合与资源相图（结论补充）}\label{app:structural-synthesis-phase-diagram}

\subsection{结构性综合结论：地址—守恒—寄存器—纤维—面积相位的核心耦合}
根据 2026-02-07 版本文档所固定的形式框架，当前体系已形成一条紧耦合的接口链：递归地址化的 $\sigma$-代数细化、以“扩展保存”将信息守恒严格化为扩展空间的单射读出、以素数寄存器/同余轴与相位纤维的紧群封装实现角色基对角化读出，并以单位圆面积相位坐标门将非紧尺度无参压缩到可审计的相位端点层。
\par
在不复述既有编号结论的前提下，本节给出一组新的\textbf{结构性综合结论}（可作为“可审计输出/可证伪接口”嵌入体系），围绕“地址—信息守恒—寄存器—纤维—单位圆面积相位”五者的核心耦合结构展开。

\subsubsection{结论 I：扩展保存下“最小记录轴”的规范四因子分解}
在“部分观测 $\Pi$ 非单射、若要求信息不丢失则须外置记录轴 $r$ 使扩展读出 $\Pi^{\mathrm e}=(\Pi,r)$ 成为单射”的制度下，“最小记录轴”并非任意对象。
当观测链同时包含：\textbf{(A)} 折叠/跨分辨率投影、\textbf{(B)} unitary slice 的强类型单参审计、\textbf{(C)} 面积相位紧化与端点读出，则任何满足扩展保存的记录轴 $r$ 都可规范分解为四个互不耦合的因子：
$$
r\ \cong\ r_{\mathrm{fin}}\times r_{\widehat{\ZZ}}\times r_{\mathrm{wind}}\times r_{\mathrm{end}}.
$$
\par
\noindent\textbf{(1) 有限缺陷因子 $r_{\mathrm{fin}}$：}对应折叠造成的有限不可逆性外置；其最小模型可由显式有限寄存器承载，使扩展读出在有限缺陷上单射化。\par
\noindent\textbf{(2) profinite 同余/素数寄存器因子 $r_{\widehat{\ZZ}}$：}对应“同余更新步/寄存器增量”在逆极限上的相容外置，并在 Pontryagin 对偶侧被角色基严格对角化为扭曲块（统一全部模数/全部角色的通道）。\par
\noindent\textbf{(3) 绕数/周期折叠因子 $r_{\mathrm{wind}}$：}对应虚方向的 $2\pi$ 周期折叠；其最小可审计实现为外置整数证书（绕数/净穿越计数），使原本折叠的角变量提升为可逆坐标，并与 Dirichlet 口径下的实现同型。\par
\noindent\textbf{(4) 端点载荷因子 $r_{\mathrm{end}}$：}对应面积相位端点压缩带来的两类不可约信息：\textbf{(i)} 离散端点分支记忆（可由有限残差承载），\textbf{(ii)} 连续精度载荷（可取 $y:=-\log|w|\in\RR_{+}$ 的平移账本承载）。二者在端点邻域严格正交分解；伸缩仅推进精度载荷的加性账本，而不改变离散端点分支类。
\par
该四因子分解将“地址缺失/$\NULL$ 的修复”从经验性修正提升为结构性分类：每一类不可逆性来源都对应一个不可约记录因子，其类型（有限、profinite、整数、端点离散$+$精度）由接口本体决定，而非实现细节决定。

\subsubsection{结论 II：离散折叠账本与面积相位精度势账本的单一混合守恒式}
体系中存在两条同构但不同域的账本恒等式：离散侧以折叠纤维的非均匀残差刻画“丢失信息”，连续侧以面积相位紧化的 Jacobian 诱导精度势刻画“尾部压缩代价”。
\par
更具体地，令 $X$ 为连续尺度变量，令 $u=\upsilon(X)$ 为其面积相位紧化坐标，并以 $h(\cdot)$ 表示微分熵，则存在严格恒等式
$$
h(X)=h(u)+\EE\bigl[P(X)\bigr],
$$
其中 $P(X)=\log\pi+\log(1+X^2)$ 为精度势（亦可视为信息账本项）。
\par
在“离散微态（折叠多重度）$\times$ 连续实轴尺度（面积相位紧化）”的状态分解下，当观测链在两坐标上分别施加 Fold 与紧化门 $\upsilon$（见第 \ref{sec:conclusion-cluster} 节的紧化坐标与精度势定义），扩展保存意义下使读出单射所需的\textbf{最小平均记录成本}（以 nats 计）满足可加的混合守恒式：
$$
\mathcal{R}_{\min}\ :=\ \kappa_m(\pi)\ +\ \EE\bigl[P(X)\bigr],
$$
其中 $\kappa_m(\pi):=\varkappa_m(\pi)$ 为折叠边缘 $\pi$ 诱导的平均纤维对数缺陷（见附录 \ref{app:auditable-profinite-chebotarev}），而 $P(X)=\log\pi+\log(1+X^2)$ 为面积相位紧化诱导的精度势。并且存在规范实现同时达到该下界：离散侧取纤维均匀提升使 KL 残差成为唯一“非均匀债务”，连续侧采用无参面积相位门使精度势成为唯一“尾部债务”。

\subsubsection{结论 III：地址完备性的失败被“精度—模积”二维配额完全刻画}
在 $\NULL$ 语义中，读出失败可分解为语义空值、协议空值与碰撞空值三类；类型层原则是：若未提供使扩展读出单射的记录轴，则读出退化为 $\NULL$。
\par
在本文接口下，碰撞空值与端点空值分别具有可审计的硬阈值结构：
\begin{itemize}
  \item \textbf{碰撞空值阈值（prime-register 配额）：}在区间读出/CRT 外置协议下，为把多对一压缩以失败概率 $\le\varepsilon$ 提升为近似可逆，prime-register 的模积预算 $P_m$ 必须超过一个显式阈值；等价地，可写为
  $$
  \log P_m\ \ge\ m\,\gamma_m(\varepsilon)+\log 2,
  $$
  其中 $\gamma_m(\varepsilon)$ 表示在分辨率 $m$ 与容许失败 $\varepsilon$ 下的可审计误差指数（由所选碰撞指标/证书协议确定）。
  \item \textbf{端点空值阈值（面积相位精度配额）：}面积相位门在有限精度下具有强误差放大律，端点邻域若精度不足以区分不同前像，则读出退化为不可核对的多值/$\NULL$；其最大可分辨动态范围与紧化步长满足
  $$
  \Delta u\ \le\ \frac{1}{\pi X_{\max}},
  $$
  其中 $\Delta u$ 为紧化坐标的量化步长，而 $X_{\max}$ 为希望在实轴上保持单值可核对读出的最大动态范围。
\end{itemize}
由此得到一个可审计的二维资源相图：离散维度由 residue/同余轴的模积配额控制，连续维度由端点精度势的量化配额控制；两者分别对应结论 I 的 $r_{\widehat{\ZZ}}$ 与 $r_{\mathrm{end}}$。

\subsubsection{结论 IV：solenoid 统一封装诱导尺度稳定地址的不变量分类}
相位纤维与 profinite 同余轴可被一维 solenoid 统一封装：
$$
\Sigma_{\mathrm{sol}}:=\bigl(\RR\times \widehat{\ZZ}_b\bigr)/\ZZ.
$$
当面积相位/Dirichlet--Cayley 扩展在端点处引入对数深度坐标并诱导伸缩半群的严格加性时，可构造伸缩不变量
$$
I:=\log\!\Bigl(\frac{\sigma\log\lambda}{|t|}\Bigr),
$$
其中 $\lambda$ 为伸缩因子，$t$ 为 unitary slice 的相位参数（例如 $\theta$ 的重参数），而 $\sigma$ 表示协议域内的固定尺度标尺。
\par
在 $\Sigma_{\mathrm{sol}}$ 的统一封装下，任何同时依赖相位与同余寄存器的可核对读出，其“尺度稳定性”由不变量类 $I$ 控制：若扩展读出所需的记录轴仅在 solenoid 的深度方向累积（即仅推进精度载荷/绕数，而不引入新的横向可见自由度），则地址等价类可按 $I$ 分层；同一层内的地址在伸缩作用下保持可寻址而不触发碰撞或协议拒绝。

\subsubsection{结论 V：地址族与角色对角化之间的傅里叶可分离性完备判据}
在地址侧，递归地址化仅执行对既有可观测事件的柱集闭包/$\sigma$-代数细化，不引入外加可测结构；在寄存器侧，profinite 轴上的提升算子在 Pontryagin 对偶侧严格分解为角色直和，扭曲通道同时对角化为扭曲矩阵块 $M_\chi$。
\par
令 $\mathcal{A}$ 为当前 PW 闭包下允许的地址集合（柱事件族），将其指示函数视为可见代数中的基元。若对任意非平凡角色 $\chi\neq 1$，存在某个地址 $a\in\mathcal{A}$ 使得该地址在 $\widehat{\ZZ}_b$ 角色基上的耦合系数不为零（等价地：$\mathcal{A}$ 在角色对偶侧对所有 $\chi\neq 1$ 具有分离能力），则协议空值的“不可见通道”被消除；此时地址完备性的剩余缺口完全退化为资源问题：碰撞配额与端点精度配额是否满足（结论 III）。
反之，若存在某个 $\chi\neq 1$ 对所有 $a\in\mathcal{A}$ 均不可见，则该 $\chi$-纤维上的谱事件无论在解析层面如何存在，都会在该闭包内不可寻址并被类型系统判为 $\NULL$。

\subsubsection{结论 VI：异常签名的谱流绕数实现与第二审计通道}
在投影词正规形与 strictification 口径下，非交换残差可压缩为常数层异常签名并进入可复核输出；另一方面，寄存器提升算子给出按角色分解的扭曲矩阵族 $M_\chi(u)$，其相位参数 $u=e^{i\theta}$ 正是 unitary slice 的可见连续轴。
\par
对每个 $\chi\neq 1$，在单位圆相位轴上考察 $M_\chi(e^{i\theta})$ 的连续谱流。将“净穿越”作为整数不变量，可取为行列式读出的绕数，例如定义
$$
\mathrm{Wind}_\chi
:=\frac{1}{2\pi i}\int_{0}^{2\pi}\frac{\dd}{\dd\theta}\log\det\!\bigl(I-M_\chi(e^{i\theta})\bigr)\,\dd\theta
\ \in\ \ZZ,
$$
其可由 Rouch{\'e}/区间证书等离线协议核验为整数（与 unitary slice 的零点证书接口一致）。
\par
若异常签名在给定切片上可被 strictification 严格吸收，则相应的谱流绕数应为零：任一围绕单位圆的闭合扫描不产生净的谱穿越计数；反之，若观测到 $\mathrm{Wind}_\chi\neq 0$，则它给出 strictification 的障碍，从而必然要求引入新的记录轴因子。在结论 I 的分类中，该新增因子只能落入 $r_{\mathrm{wind}}$ 或 $r_{\mathrm{end}}$ 的加厚，而不可能由“末端追加少量位”的局部修正实现。

\subsubsection{收束性总结：五者统一的可审计输出形态}
上述六条结论将五者（地址/守恒/寄存器/纤维/面积相位）的耦合压缩为可审计、可计算、可证伪的结构性陈述：
\begin{enumerate}
  \item 地址完备性可在对偶侧转写为“覆盖全部扭曲通道”的傅里叶分离性质（结论 V）。
  \item 信息守恒在体系内等价于“存在可分解的记录轴，使扩展读出单射”，且最小记录轴具有规范四因子分解（结论 I）。
  \item 面积相位门诱导的精度势与折叠账本同型，二者可合并为单一混合守恒式并给出最小平均记录成本（结论 II）。
  \item solenoid 统一封装与端点深度账本给出尺度稳定地址的不变量分类，从而把“推进伸缩链是否必触发 $\NULL$”化为可审计结构问题（结论 IV）。
  \item 异常签名与谱读出形成双通道审计闭环：以谱流绕数作为独立的整数审计物，与 strictification/协议拒绝互为交叉核验（结论 VI）。
\end{enumerate}

\subsection{扩展保存下的“可寻址性”判据：规范差相图与记录轴信息率}
根据 2026-02-07 版本的接口链条（递归地址化的 $\sigma$-代数细化、扩展保存、素数寄存器/同余轴与相位纤维的紧群封装、以及单位圆面积相位门的端点压缩），可将“可寻址且非 $\NULL$”的失败机制从抽象完备性（completeness）缺口推进为一组可核对的统计/协议判据。以下给出与正文既有编号结果互补的一组结构性结论，聚焦于“地址缺失/记录轴不足”的统一刻画及其资源相图。

\subsubsection{扩展保存下的可寻址性判据：地址缺失与记录轴不足的统一刻画}
在部分观测语义中，观测映射一般多对一。为严格表达“信息不丢失”，引入记录映射并要求扩展观测单射。
\par
设 $\Pi:\Omega\to X$ 为观测映射（例如固定分辨率 $m$ 的折叠 $\Fold_m:\Omega_m\twoheadrightarrow X_m$）。称 $r:\Omega\to R$ 为一个记录轴，若扩展观测
$$
\Pi_{\mathrm e}(\omega):=\bigl(\Pi(\omega),r(\omega)\bigr)\in X\times R
$$
为单射。此时“扩展保存”即为：在扩展空间 $X\times R$ 上存在可逆回放器恢复 $\omega$。
\par
在“地址先于取值”的类型约束下，若未提供使 $\Pi_{\mathrm e}$ 单射的记录轴，则读出退化为 $\NULL$；并且该 $\NULL$ 可分解为三类机制：
\begin{itemize}
  \item \textbf{语义空值（地址缺失）：}对象在当前地址闭包下不可呈现，因此不存在可核对读出；
  \item \textbf{协议空值（类型系统禁止）：}输入不属于被允许的强类型接口域（例如 unitary slice 仅允许 $(L,\theta,p)$），从而被拒绝；
  \item \textbf{碰撞空值（侧信息预算不足）：}即使语义允许且协议放行，但记录轴容量不足以将多对一观测提升为近似单射，导致不可唯一重建。
\end{itemize}
该三分解把“地址完备性失败”的抽象缺口细化为可审计的失败类型：前两类为类型层/闭包层缺口，第三类为资源不足的统计缺口。

\subsubsection{偏置输入下的规范差相图：从窗口不相容到记录轴速率曲线}
在跨分辨率降阶时，朴素截断位与折叠感知 restriction 位一般不相容。令 $G_m(\omega)$ 为分辨率 $m$ 下两种读出在 $m$ 位窗口上的逐位失配总数（规范差），即
$$
G_m(\omega):=\sum_{i=1}^{m}\mathbf{1}\{X_i(\omega)\neq Y_i(\omega)\}.
$$
在 IID Bernoulli($p$) 基线下（$\PP(X_i=1)=p$），规范差的指数矩母函数可由一个有限维加权邻接矩阵表示，进而得到可审计的压力接口。

\begin{theorem}[Bernoulli($p$) 基线的规范差压力矩阵]\label{thm:conclusion-normdiff-pressure-matrix}
对任意 $p\in(0,1)$，定义 $4\times 4$ 矩阵族
$$
A_\theta(p)=
\begin{pmatrix}
1-p & p e^{\theta} & 0 & p\\
0 & 0 & p e^{\theta} & 0\\
(1-p)e^{\theta} & 1 & 0 & 0\\
1-p & 0 & 0 & 0
\end{pmatrix},
$$
并记谱半径为 $\rho(\cdot)$。定义规范差压力函数
$$
P_G(\theta;p):=\log\rho\!\bigl(A_\theta(p)\bigr).
$$
则 $P_G(\theta;p)$ 给出规范差的可审计压力接口：其一阶/二阶/三阶导数分别对应极限均值、方差密度与三阶累积密度（见下述定理）。
\end{theorem}

\begin{theorem}[规范差密度曲线的闭式表达]\label{thm:conclusion-normdiff-density-closedform}
令
$$
g(p):=\lim_{m\to\infty}\frac{\EE[G_m]}{m}.
$$
则有闭式
$$
g(p)=\frac{p^{2}(3-2p)}{1+p^{3}}.
$$
并且 $g$ 在 $p\in(0,1)$ 上存在唯一内点极大值点 $p_\star$，其由方程
$$
p_\star^{3}+2p_\star-2=0
$$
确定。
\end{theorem}

\begin{theorem}[内点二元联合律的 $p$-闭式化]\label{thm:conclusion-normdiff-jointlaw}
在内点极限口径下，朴素截断位 $X\in\{0,1\}$ 与折叠感知位 $Y\in\{0,1\}$ 的联合分布为
$$
\PP(X=0,Y=0)=\frac{1-p-p^{2}+2p^{3}-p^{4}}{1+p^{3}},\qquad
\PP(X=0,Y=1)=\frac{p^{2}(1-p)}{1+p^{3}},
$$
$$
\PP(X=1,Y=0)=\frac{p^{2}(2-p)}{1+p^{3}},\qquad
\PP(X=1,Y=1)=\frac{p\bigl(p^{3}+(1-p)^{2}\bigr)}{1+p^{3}}.
$$
因此边缘律严格保留输入偏置 $\PP(X=1)=p$，并且单点失配概率满足
$$
\PP(X\neq Y)=g(p).
$$
\end{theorem}

\subsubsection{规范差的涨落常数族：CLT 方差密度与三阶累积密度的 $p$-闭式}
压力函数 $\theta\mapsto P_G(\theta;p)$ 的高阶导数给出规范差的涨落指纹族；其中二阶导数对应 CLT 方差密度，三阶导数对应最低阶“非自对偶”偏斜指纹。

\begin{theorem}[CLT 方差密度]\label{thm:conclusion-normdiff-clt-variance}
存在方差密度 $\sigma^{2}(p)>0$ 使得
$$
\frac{G_m-g(p)m}{\sqrt m}\Longrightarrow \mathcal N\bigl(0,\sigma^{2}(p)\bigr),
$$
并且
$$
\sigma^{2}(p)=
\frac{p^{2}\bigl(-21p^{6}+27p^{5}-20p^{4}+50p^{3}-43p^{2}-2p+9\bigr)}{(1+p^{3})^{3}}.
$$
\end{theorem}

\begin{theorem}[三阶累积密度]\label{thm:conclusion-normdiff-third-cumulant}
记 $P_G^{(3)}(0;p)$ 为 $\theta\mapsto P_G(\theta;p)$ 在 $\theta=0$ 处的三阶导数，则
$$
P_G^{(3)}(0;p)=
\frac{p^{2}\bigl(117p^{12}-567p^{11}+692p^{10}-1602p^{9}+3897p^{8}-3370p^{7}+1134p^{6}-1377p^{5}+1264p^{4}+156p^{3}-381p^{2}+10p+27\bigr)}{(1+p^{3})^{5}}.
$$
\end{theorem}

\subsubsection{从规范差到“记录轴最小信息率”：条件熵的闭式接口}
由扩展保存，若希望把多对一观测提升为单射或近似单射，则必须外置记录轴。对于二元极限口径，可将“每步所需记录轴的最小平均信息率”定义为条件熵
$$
H(X\mid Y):=-\sum_{x,y\in\{0,1\}}\PP(X=x,Y=y)\log \PP(X=x\mid Y=y),
$$
它刻画：仅观察 $Y$ 而要无损恢复 $X$ 时，记录轴的最小平均信息开销。由定理 \ref{thm:conclusion-normdiff-jointlaw} 的联合律，$H(X\mid Y)$ 在 Bernoulli($p$) 基线下成为显式可计算函数；当记录轴预算不足以覆盖该信息率时，即发生“碰撞空值”型 $\NULL$。

\subsubsection{与单位圆面积相位门的统一：连续端点精度记录轴与离散规范差同构}
面积相位门由
$$
\vartheta(x):=2\arctan(x),\qquad \iota(x):=e^{i\vartheta(x)}=\frac{1+ix}{1-ix}\in\TT
$$
给出：它把无穷远压缩到单位圆边界点 $-1$，并将 Cauchy 权重严格推前为单位圆 Haar 测度。更精确地，有换元恒等式
$$
\frac{1}{\pi}\frac{\dd x}{1+x^2}=\frac{1}{2\pi}\,\dd\vartheta,
$$
从而 $\iota$ 将 $\frac{1}{\pi}(1+x^2)^{-1}\dd x$ 推前为 $\TT$ 上的均匀测度。
\par
在该门语言下，实轴上的“相位相加”在射影紧化上诱导交换群律
$$
x\oplus y=\frac{x+y}{1-xy},
$$
并满足群同态 $\iota(x\oplus y)=\iota(x)\iota(y)$，从而把相位纤维的可逆运算实现为单位圆乘法。
\par
更关键的是：紧化把动态范围压入末位精度。若取紧化坐标 $u=\upsilon(x)=(\pi^{-1})\arctan(x)$，则
$$
\frac{\dd x}{\dd u}=\pi(1+x^{2}),
$$
因此当 $|x|$ 增大时，读出精度必须通过额外末位记录（更细的 $u$-分辨率）来补偿；这与离散折叠中的“规范差必须外置为记录轴”属于同一扩展保存范式。
\par
最后，离散侧的折叠门给出信息账本恒等式（定理 \ref{thm:ledger-identity}），把可见熵、纤维体积与残差信息严格分解；连续侧的面积相位门给出微分熵恒等式 $h(X)=h(u)+\EE[P(X)]$，将无穷端点转换为精度势/末位预算的可审计接口。两者共同提供了一个统一的“记录轴预算”语言：离散侧的纤维缺陷与连续侧的精度势在同一账本上以可加方式出现，从而把 $\NULL$ 的资源机制压缩为可核验的预算不等式。

